% !TeX root = ../教材题库_全部条目_转写工作稿.tex
% Source: XJ-B1-C03-RV-P090-Y0720
% 第3章 / 章末复习 复习题三 / 教材 p.90 / PDF 第97页
% OCR未校对：公式与文字均可能有误；未完成数学TeX转写；禁止直接发给学生
% 原题及相邻原解截图（核对用）：../题目原图/XJ-B1-C03-RV-P090-Y0720-01.webp
% 本文件只在转写完成后作为题干使用；OCR 条目尚不是可靠题干。
\begin{bankitem}{复习题三 第4题}{第3章 / 章末复习 复习题三 / 教材 p.90 / PDF 第97页}{XJ-B1-C03-RV-P090-Y0720}{OCR未校对：公式与文字均可能有误；未完成数学TeX转写；禁止直接发给学生}
\begin{ocrdraft}
[教材p.90 / PDF第97页]
\par
4．已知函数y一丆0）的对应关系如下表，函数y=g(x)的图象是如图所示的
\par
曲线ABC，其中A(1，3），B(2，1），C（3，2），则丆（g（2））的值为
\par
3
\par
2
\par
1
\par
2
\par
3
\par
2
\par
3
\par
0
\par
2
\par
（第4题）
\end{ocrdraft}
\end{bankitem}
